\input{Datos.txt}
%%GIRAR TODO EL DOCUMENTO
%\documentclass[a4paper,landscape]{article}
\documentclass[letterpaper,11pt,landscape]{article}
%GIRAR UN ESPACIO ESPECIFICO 
\usepackage{pdflscape}
\usepackage{array}

\usepackage{graphicx}
\usepackage{longtable}
\usepackage{booktabs}

\usepackage{datatool}
%\DTLloaddb{ALTAS}{ALTAS.csv}

\DTLloaddb{ALTASINTER}{ALTASINTER.csv}
\DTLloaddb{ALTASDES}{ALTASDES.csv}
\DTLloaddb{ALTASUSUA}{ALTASUSUA.csv}
\DTLloaddb{ALTASADMIN}{ALTASADMIN.csv}
\DTLloaddb{ALTASOTRO}{ALTASOTRO.csv}

%\DTLloaddb{BAJAS}{BAJAS.csv}

\DTLloaddb{BAJASINTER}{BAJASINTER.csv}
\DTLloaddb{BAJASDES}{BAJASDES.csv}
\DTLloaddb{BAJASUSUA}{BAJASUSUA.csv}
\DTLloaddb{BAJASADMIN}{BAJASADMIN.csv}
\DTLloaddb{BAJASOTRO}{BAJASOTRO.csv}
%\DTLloaddb{CAMBIOS}{CAMBIOS.csv}
\usepackage{ifthen}
\usepackage{filecontents}
\usepackage{setspace}
\usepackage{tabularx}


% *** Idioma ***
\usepackage[spanish,es-nodecimaldot]{babel}
\usepackage[utf8]{inputenc}

% *** Fuente ***
%\usepackage{lmodern}
%\usepackage[T1]{fontenc}
%%FUENTE NOTO
\usepackage{fontspec}
\setmainfont{Noto Sans}

% *** Colores ***
\usepackage{xcolor}
\definecolor{azul}{RGB}{47,125,188} 
\definecolor{dorado}{RGB}{190,150,91} 
\definecolor{gris}{rgb}{0.752,0.752,0.752}
\definecolor{morado}{RGB}{28,28,129}
\definecolor{gristabla}{rgb}{0.85,0.85,0.85}
\definecolor{BarleyWhite}{rgb}{1,0.949,0.8}
\definecolor{Spindle}{rgb}{0.741,0.843,0.933}
\definecolor{Zanah}{rgb}{0.886,0.941,0.85}
\definecolor{black}{rgb}{0,0,0}


% *** Margenes y Tamaño de Hoja
\usepackage{geometry}
\geometry{
left = 15mm,
right = 15mm,
top = 45mm,
bottom = 15mm,
headheight = 28mm
}

% *** Imagenes ***
\usepackage{graphicx}
\graphicspath{./imagenes}
\usepackage{float}
\usepackage{subcaption}
\usepackage{wrapfig}
\usepackage{tabularray}

% *** Fecha ***
\usepackage[useregional]{datetime2}

% *** Encabezado ***
\usepackage{lastpage} 
\usepackage{fancyhdr}
\pagestyle{fancy}
\renewcommand{\headrulewidth}{0pt}
\lhead{
		\begin{tblr}{
		width = \linewidth,
		colspec = {Q[230]Q[507]Q[100]Q[100]},
		row{2} = {c},
		row{3} = {c},
		column{2} = {c},
		cell{1}{1} = {r=4}{c},
		cell{1}{2} = {r=3}{},
		cell{1}{3} = {c},
		cell{1}{4} = {c},
                     cell{4}{3} = {c},
                     cell{4}{4} = {c},
		%cell{4}{3} = {c=2}{c},
		vlines,
		hline{1,5} = {-}{},
		hline{2-3} = {3-4}{},
		hline{4} = {2-4}{},
	}
	\includegraphics[width=0.2\textwidth]{imagenes/Logo_de_la_Conagua} & {\textbf {COMISIÓN NACIONAL DEL AGUA}\\Gerencia de Tecnologías de Información y Comunicaciones. \\ Subgerencia de Soporte Técnico, Telecomunicaciones y Seguridad.} & \footnotesize HOJA    & {\footnotesize \thepage\ de \pageref{LastPage}} \\
	&                                                                                       &\footnotesize PROCESO   & \footnotesize ACMB \\
	&                                                                                       &{ \footnotesize FECHA \\ 07-04-2025} & {\footnotesize Último Cambio\\ Liberado }    \\
	&\textbf { Solicitud de alta baja o cambio en la infraestructura de seguridad de la Comisión Nacional del Agua.}                                       &  { \footnotesize No. Formato \\ \NOFORMATO}       & { \footnotesize ANEXO 17 FORMATO 1}
\end{tblr}
}
\chead{}
\rhead{}
% *** Pie de Pagina
%\renewcommand{\footrulewidth}{0.5pt}
\lfoot{}
\cfoot{}
\rfoot{}


%No cortar palabras
\usepackage[none]{hyphenat}
\begin{document}
\sloppy
%%%INFORMACIÓN DEL SOLICITANTE Y ENLACE INFORMATICO 
\newboolean{soli}
\setboolean{soli}{\SOLI}
\ifthenelse{\boolean{soli}}{
\vspace{-30pt}
\begin{longtblr}[
	label = none,
	entry = none,
	]{
		width = \linewidth,
		colspec = {Q[100]Q[300]Q[90]Q[90]},
                     row{1}={dorado},                    
		row{1}={c},
                     cell{1}{1} = {c=4}{0.874\linewidth},	
                     cell{2}{3}={r},
                     cell{3}{2} = {c=3}{0.679\linewidth},                    
                     cell{4}{3}={r},	
                     cell{4}{2} = {c=3}{0.679\linewidth},                     
		hlines,
		vlines,
	}
        
   \textbf{Información del solicitante:} &  &      &  \\
             Nombre:                &    \NOMS &   {Teléfono/Extensión:}&\EXTS\\
             Puesto:&\PUESTOS\\  Área: &\AREAS            
\end{longtblr}}

\newboolean{enlace}
\setboolean{enlace}{\ENLACE}
\ifthenelse{\boolean{enlace}}{
\vspace{-40pt}
\begin{longtblr}[
	label = none,
	entry = none,
	]{
		width = \linewidth,
		colspec = {Q[100]Q[300]Q[90]Q[90]},
                     row{1}={dorado},                    
		row{1}={c},
                     cell{1}{1} = {c=4}{0.874\linewidth},	
                     cell{2}{3}={r},                     
		hlines,
		vlines,
	}
             \textbf{Información del enlace informático:} &  &      &  \\
             Nombre:                &   \NOMEI  &   {Teléfono/Extensión:}&\EXTEI       
           
\end{longtblr}}





\vspace{-15pt}
\begin{longtblr}[
	label = none,
	entry = none,
	]{
		width = \linewidth,
		colspec = {Q[203]Q[200]Q[200]Q[200]},
                     row{1}={dorado},
		row{1}={c},
                     row{5}={dorado},
		row{5}={c},
                     cell{1}{1} = {c=4}{0.679\linewidth},
		cell{2}{3} = {r},
		cell{2}{4} = {c},
                     cell{5}{1} = {c=4}{0.679\linewidth},
                     cell{7}{3} = {c=2}{0.379\linewidth},
                     cell{3}{2} = {c=3}{0.679\linewidth},
                     cell{4}{2} = {c=3}{0.679\linewidth},
                     %cell{5}{2} = {c=3}{0.679\linewidth},
                     row{6}={c},
                     cell{7}{1}={c},
                     cell{7}{2}={c},                     
                     	hlines,
		vlines,
	}
          
          \textbf{Información para el registro de solicitud:} &  &      &  \\
	Numero de ticket: & \NOTICKET & {Fecha de solicitud:}   & \today       \\
          {Memorando / Atenta nota:}&\MEMO\\
         Solicitud:&  \DESCBREVE\\
            \textbf {Tipo de Cambio:} &\\
            Intersistemas (\INTER) & Administrativo (\ADMIN) & Desarrollador (\DES )& Usuario(\USUA)\\
           Otro (\OTRO )& Describa brevemente:&\DESOTRO\\   
                 
\end{longtblr}



\vspace{-15pt}
\begin{longtblr}[
	label = none,
	entry = none,
	]{
		width = \linewidth,
		colspec = {Q[935]},
		cells = {c},
                     row{1}={dorado},
		row{2} = {},
		hlines,
		vlines,
	}
	\textbf{Descripción detallada} \\ \DESDET
\end{longtblr}


%%%TABLAS DE ALTAS BAJAS Y CAMBIOS
%%%PARA ALTAS
\clearpage
\newboolean{altasinter}
\setboolean{altasinter}{\ALTASINTER} 
\newboolean{altasadmin}
\setboolean{altasadmin}{\ALTASADMIN} 
\newboolean{altasdes}
\setboolean{altasdes}{\ALTASDES} 
\newboolean{altasusua}
\setboolean{altasusua}{\ALTASUSUA} 

%%PARA BAJAS
\newboolean{bajasinter}
\setboolean{bajasinter}{\BAJASINTER} 
\newboolean{bajasadmin}
\setboolean{bajasadmin}{\BAJASADMIN} 
\newboolean{bajasdes}
\setboolean{bajasdes}{\BAJASDES} 
\newboolean{bajasusua}
\setboolean{bajasusua}{\BAJASUSUA} 





%%%%%%%
%%%%INTERSISTEMAS
\newboolean{inter}
\setboolean{inter}{\INTERBOOL} 
\ifthenelse{\boolean{inter}}
{
%%ALTASINTERSISTEMAS

\ifthenelse{\boolean{altasinter}}
{
\vspace{-25pt}
\begin{longtblr}[
	label = none,
	entry = none,
	]{
		width = \linewidth,
		colspec = {Q[935]},
		cells = {c},
                     row{1}={dorado},                     
		%hlines,
		%vlines,
	}
	\textbf{ALTAS} \\Sección A.- Comunicación intersistemas
\end{longtblr}
%%%PARA PODER ALMACENAR TODAS LAS FILAS DEL DOCUMENTOS CSV ES NESARIO TODO ESTO
\vspace{-30pt}
 \begin{longtblr}[
 label = none,
 entry = none,
 ]{
  width = \linewidth,
  colspec = {Q[2]Q[4]Q[6]Q[12]Q[4]Q[6]Q[12]Q[5]Q[4]},                     
  hlines,
  vlines,
                     cells={c},
 }
\textbf{No°}&\textbf {Sistema\\ Origen}&\textbf{Función o rol de \\dipositivo(s) \\origen}&
\textbf{IP/NAT dispositivo(s) \\origen/NAT}&\textbf{Sistema\\ Destino} &
\textbf{Función o rol de \\dispositivo(s) \\destino}&\textbf{IP/NAT dispositivo(s) \\destino/NAT}&
\textbf{Protocolo\\ TCP o UDP}&\textbf{Puertos}
\end{longtblr}
%%%%INTERSISTEMAS
\DTLforeach*{ALTASINTER}{\No=N, \SistemaOri=SO, \FuncionOri=FRO, \IPOri=IPO, \SistemaDes=SD, \FuncionDes=FRD, \IPDes=IPD, \Protocolo=PRO, \Puertos=PUER}
{
\vspace{-37pt}
 \begin{longtblr}[
 label = none,
 entry = none,
 ]{
  width = \linewidth,
  colspec = {Q[2]Q[4]Q[6]Q[12]Q[4]Q[6]Q[12]Q[5]Q[4]},                     
  hlines,
  vlines,
                     cells={c},
 }
\No&\SistemaOri&\FuncionOri&\IPOri&\SistemaDes& \FuncionDes&\IPDes&\Protocolo& \Puertos
\end{longtblr}
}
%%%%%TEMPORALIDAD INTERSISTEMAS 
\vspace{-20pt}
\begin{longtblr}[
	label = none,
	entry = none,
	]{
		width = \linewidth,
		colspec = {Q[935]},
		cells = {c},
                     row{1}={dorado},
		row{2} = {},
		hlines,
		vlines,
	}
	\textbf{Temporalidad de Intersistemas} \\ \TEMPOINTER
\end{longtblr}
\clearpage
}%%FIN CONDICIONAL IF

%%%%%%%%
%%BAJAS INTERSISTEMAS

\ifthenelse{\boolean{bajasinter}}
{
\vspace{-25pt}
\begin{longtblr}[
	label = none,
	entry = none,
	]{
		width = \linewidth,
		colspec = {Q[935]},
		cells = {c},
                     row{1}={dorado},
		%hlines,
		%vlines,
	}
	\textbf{BAJAS} \\Sección A.- Comunicación intersistemas
\end{longtblr}
%%BAJASintersistemas
\vspace{-30pt}
 \begin{longtblr}[
 label = none,
 entry = none,
 ]{
  width = \linewidth,
  colspec = {Q[2]Q[4]Q[6]Q[12]Q[4]Q[6]Q[12]Q[5]Q[4]},                     
  hlines,
 vlines,
                     cells={c},
 }
\textbf{No°}&\textbf {Sistema\\ Origen}&\textbf{Función o rol de \\dipositivo(s) \\origen}&
\textbf{IP/NAT dispositivo(s) \\origen/NAT}&\textbf{Sistema\\ Destino} &
\textbf{Función o rol de \\dispositivo(s) \\destino}&\textbf{IP/NAT dispositivo(s) \\destino/NAT}&
\textbf{Protocolo\\ TCP o UDP}&\textbf{Puertos}
\end{longtblr}
%%%%INTERSISTEMAS
\DTLforeach*{BAJASINTER}{\No=N, \SistemaOri=SO, \FuncionOri=FRO, \IPOri=IPO, \SistemaDes=SD, \FuncionDes=FRD, \IPDes=IPD, \Protocolo=PRO, \Puertos=PUER}
{
\vspace{-37pt}
 \begin{longtblr}[
 label = none,
 entry = none,
 ]{
  width = \linewidth,
  colspec = {Q[2]Q[4]Q[6]Q[12]Q[4]Q[6]Q[12]Q[5]Q[4]},                     
  hlines,
 vlines,
                     cells={c},
 }
\No&\SistemaOri&\FuncionOri&\IPOri&\SistemaDes& \FuncionDes&\IPDes&\Protocolo& \Puertos
\end{longtblr}
}
}%%FIN CONDICIONAL
}%%FIN DE INTERSISTEMAS


%%%%ADMINISTRATIVOS
\newboolean{admin}
\setboolean{admin}{\ADMINBOOL} 
\ifthenelse{\boolean{admin}}
{
%%ALTASADMIN

\ifthenelse{\boolean{altasadmin}}
{
\vspace{-25pt}
\begin{longtblr}[
	label = none,
	entry = none,
	]{
		width = \linewidth,
		colspec = {Q[935]},
		cells = {c},
                     row{1}={dorado},
		%hlines,
		%vlines,
	}
	\textbf{ALTAS} \\Sección B.- Accesos administrativos
\end{longtblr}
%%%PARA PODER ALMACENAR TODAS LAS FILAS DEL DOCUMENTOS CSV ES NESARIO TODO ESTO
\vspace{-30pt}
 \begin{longtblr}[
 label = none,
 entry = none,
 ]{
  width = \linewidth,
  colspec = {Q[2]Q[7]Q[12]Q[10]Q[8]Q[15]Q[5]Q[5]},                     
  hlines,
 vlines,
                     cells={c},
 }
\textbf{No°}&\textbf {Nombre(s) de \\ administrador\\(es)}&\textbf{IP anfitrión(es) \\origen}&
\textbf{Sistema asociado a \\ anfitrión(es) destino}&
\textbf{Función o rol de \\anfitrión(es) \\destino}&\textbf{IP/NAT anfitrión(es) \\destino/NAT}&
\textbf{Protocolo\\ TCP o UDP}&\textbf{Puertos}
\end{longtblr}
%%%%ADMINISTRATIVOS
\DTLforeach*{ALTASADMIN}{\No=N, \NombreAdmin=SO, \IPOri=IPO, \SistemaDes=SD, \FuncionDes=FRD, \IPDes=IPD, \Protocolo=PRO, \Puertos=PUER}
{
\vspace{-37pt}
 \begin{longtblr}[
 label = none,
 entry = none,
 ]{
  width = \linewidth,
  colspec ={Q[2]Q[7]Q[12]Q[10]Q[8]Q[15]Q[5]Q[5]},                     
  hlines,
 vlines,
                     cells={c},
 }
\No&\NombreAdmin&\IPOri&\SistemaDes& \FuncionDes&\IPDes&\Protocolo& \Puertos
\end{longtblr}
}
%%%%%TEMPORALIDAD ADMINISTRATIVOS 
\vspace{-20pt}
\begin{longtblr}[
	label = none,
	entry = none,
	]{
		width = \linewidth,
		colspec = {Q[935]},
		cells = {c},
                     row{1}={dorado},
		row{2} = {},
		hlines,
		vlines,
	}
	\textbf{Temporalidad de Administrativos} \\ \TEMPOADMIN
\end{longtblr}
}
%%%%%%%%

%%BAJASADMIN
\ifthenelse{\boolean{bajasadmin}}
{
\vspace{-25pt}
\begin{longtblr}[
	label = none,
	entry = none,
	]{
		width = \linewidth,
		colspec = {Q[935]},
		cells = {c},
                     row{1}={dorado},
		%hlines,
		%vlines,
	}
	\textbf{BAJAS} \\Sección B.- Accesos administrativos
\end{longtblr}
%%%PARA PODER ALMACENAR TODAS LAS FILAS DEL DOCUMENTOS CSV ES NESARIO TODO ESTO
\vspace{-30pt}
 \begin{longtblr}[
 label = none,
 entry = none,
 ]{
  width = \linewidth,
  colspec = {Q[2]Q[7]Q[12]Q[10]Q[8]Q[15]Q[5]Q[5]},                     
  hlines,
 vlines,
                     cells={c},
 }
\textbf{No°}&\textbf {Nombre(s) de \\ administrador\\(es)}&\textbf{IP anfitrión(es) \\origen}&
\textbf{Sistema asociado a \\ anfitrión(es) destino}&
\textbf{Función o rol de \\anfitrión(es) \\destino}&\textbf{IP/NAT anfitrión(es) \\destino/NAT}&
\textbf{Protocolo\\ TCP o UDP}&\textbf{Puertos}
\end{longtblr}
%%%%ADMINISTRATIVOS
\DTLforeach*{BAJASADMIN}{\No=N, \NombreAdmin=SO, \IPOri=IPO, \SistemaDes=SD, \FuncionDes=FRD, \IPDes=IPD, \Protocolo=PRO, \Puertos=PUER}
{
\vspace{-37pt}
 \begin{longtblr}[
 label = none,
 entry = none,
 ]{
  width = \linewidth,
  colspec ={Q[2]Q[7]Q[12]Q[10]Q[8]Q[15]Q[5]Q[5]},                     
  hlines,
 vlines,
                     cells={c},
 }
\No&\NombreAdmin&\IPOri&\SistemaDes& \FuncionDes&\IPDes&\Protocolo& \Puertos
\end{longtblr}
}
}
}%%FIN DE ADMINISTRATIVOS


%%%%DESARROLLADORES
\newboolean{des}
\setboolean{des}{\DESBOOL} 
\ifthenelse{\boolean{des}}
{
%%ALTASDES

\ifthenelse{\boolean{altasdes}}
{
\vspace{-25pt}
\begin{longtblr}[
	label = none,
	entry = none,
	]{
		width = \linewidth,
		colspec = {Q[935]},
		cells = {c},
                     row{1}={dorado},
		%hlines,
		%vlines,
	}
	\textbf{ALTAS} \\Sección C.- Accesos de desarrollador
\end{longtblr}
%%%PARA PODER ALMACENAR TODAS LAS FILAS DEL DOCUMENTOS CSV ES NECESARIO TODO ESTO
\vspace{-30pt}
 \begin{longtblr}[
 label = none,
 entry = none,
 ]{
  width = \linewidth,
  colspec = {Q[2]Q[7]Q[12]Q[10]Q[8]Q[15]Q[5]Q[5]},                     
  hlines,
 vlines,
                     cells={c},
 }
\textbf{No°}&\textbf {Nombre(s) de \\ desarrollador\\(es)}&\textbf{IP anfitrión(es) \\origen}&
\textbf{Sistema asociado a \\ anfitrión(es) destino}&
\textbf{Función o rol de \\anfitrión(es) \\destino}&\textbf{IP/NAT anfitrión(es) \\destino/NAT}&
\textbf{Protocolo\\ TCP o UDP}&\textbf{Puertos}
\end{longtblr}
%%%%DESARROLLADORES
\DTLforeach*{ALTASDES}{\No=N, \NombreDes=SO, \IPOri=IPO, \SistemaDes=SD, \FuncionDes=FRD, \IPDes=IPD, \Protocolo=PRO, \Puertos=PUER}
{
\vspace{-37pt}
 \begin{longtblr}[
 label = none,
 entry = none,
 ]{
  width = \linewidth,
  colspec ={Q[2]Q[7]Q[12]Q[10]Q[8]Q[15]Q[5]Q[5]},                     
  hlines,
vlines,
                     cells={c},
 }
\No&\NombreDes&\IPOri&\SistemaDes& \FuncionDes&\IPDes&\Protocolo& \Puertos
\end{longtblr}
}
%%%%%TEMPORALIDAD DESARROLLADORES 
\vspace{-20pt}
\begin{longtblr}[
	label = none,
	entry = none,
	]{
		width = \linewidth,
		colspec = {Q[935]},
		cells = {c},
                     row{1}={dorado},
		row{2} = {},
		hlines,
		vlines,
	}
	\textbf{Temporalidad de Desarrolladores} \\ \TEMPODES
\end{longtblr}
}
%%%%%%%%

%%BAJAS DES
\ifthenelse{\boolean{bajasdes}}
{
\vspace{-25pt}
\begin{longtblr}[
	label = none,
	entry = none,
	]{
		width = \linewidth,
		colspec = {Q[935]},
		cells = {c},
                     row{1}={dorado},
		%hlines,
		%vlines,
	}
	\textbf{BAJAS} \\Sección C.- Accesos de desarrollador
\end{longtblr}
%%%PARA PODER ALMACENAR TODAS LAS FILAS DEL DOCUMENTOS CSV ES NECESARIO TODO ESTO
\vspace{-30pt}
 \begin{longtblr}[
 label = none,
 entry = none,
 ]{
  width = \linewidth,
  colspec = {Q[2]Q[7]Q[12]Q[10]Q[8]Q[15]Q[5]Q[5]},                     
  hlines,
  vlines,
                     cells={c},
 }
\textbf{No°}&\textbf {Nombre(s) de \\ desarrollador\\(es)}&\textbf{IP anfitrión(es) \\origen}&
\textbf{Sistema asociado a \\ anfitrión(es) destino}&
\textbf{Función o rol de \\anfitrión(es) \\destino}&\textbf{IP/NAT anfitrión(es) \\destino/NAT}&
\textbf{Protocolo\\ TCP o UDP}&\textbf{Puertos}
\end{longtblr}
%%%%DESARROLLADORES
\DTLforeach*{BAJASDES}{\No=N, \NombreDes=SO, \IPOri=IPO, \SistemaDes=SD, \FuncionDes=FRD, \IPDes=IPD, \Protocolo=PRO, \Puertos=PUER}
{
\vspace{-37pt}
 \begin{longtblr}[
 label = none,
 entry = none,
 ]{
  width = \linewidth,
  colspec ={Q[2]Q[7]Q[12]Q[10]Q[8]Q[15]Q[5]Q[5]},                     
  hlines,
 vlines,
                     cells={c},
 }
\No&\NombreDes&\IPOri&\SistemaDes& \FuncionDes&\IPDes&\Protocolo& \Puertos
\end{longtblr}
}
}
}%%FIN DE DESARROLLADORES


%%%%USUARIO
\newboolean{usua}
\setboolean{usua}{\USUABOOL} 
\ifthenelse{\boolean{usua}}
{
%%ALTAS USUA

\ifthenelse{\boolean{altasusua}}
{
\vspace{-25pt}
\begin{longtblr}[
	label = none,
	entry = none,
	]{
		width = \linewidth,
		colspec = {Q[935]},
		cells = {c},
                     row{1}={dorado},
		%hlines,
		%vlines,
	}
	\textbf{ALTAS} \\Sección D.- Accesos de usuario
\end{longtblr}
%%%PARA PODER ALMACENAR TODAS LAS FILAS DEL DOCUMENTOS CSV ES NECESARIO TODO ESTO
\vspace{-30pt}
 \begin{longtblr}[
 label = none,
 entry = none,
 ]{
  width = \linewidth,
  colspec = {Q[2]Q[10]Q[12]Q[10]Q[8]Q[15]Q[5]Q[5]},                     
  hlines,
 vlines,
                     cells={c},
 }
\textbf{No°}&\textbf {Nombre(s) de \\ usuario(s)\\zona(s)\\origen}&\textbf{Segmento(s)/IP(s) \\origen}&
\textbf{Sistema asociado a \\ anfitrión(es) destino}&
\textbf{Rol de \\anfitrión(es) \\destino}&\textbf{IP/NAT anfitrión(es) \\destino/NAT}&
\textbf{Protocolo\\ TCP o UDP}&\textbf{Puertos}
\end{longtblr}
%%%%USUARIO
\DTLforeach*{ALTASUSUA}{\No=N, \NombreUsua=SO, \IPOri=IPO, \SistemaDes=SD, \FuncionDes=FRD, \IPDes=IPD, \Protocolo=PRO, \Puertos=PUER}
{
\vspace{-37pt}
 \begin{longtblr}[
 label = none,
 entry = none,
 ]{
  width = \linewidth,
  colspec ={Q[2]Q[10]Q[12]Q[10]Q[8]Q[15]Q[5]Q[5]},                     
  hlines,
vlines,
                     cells={c},
 }
\No&\NombreUsua&\IPOri&\SistemaDes& \FuncionDes&\IPDes&\Protocolo& \Puertos
\end{longtblr}
}
%%%%%TEMPORALIDAD USUARIOS 
\vspace{-20pt}
\begin{longtblr}[
	label = none,
	entry = none,
	]{
		width = \linewidth,
		colspec = {Q[935]},
		cells = {c},
                     row{1}={dorado},
		row{2} = {},
		hlines,
		vlines,
	}
	\textbf{Temporalidad de Usuarios} \\ \TEMPOUSUA
\end{longtblr}%%fin de tempo usua
\clearpage
}
%%%%%%%%

%%BAJASUSUA
\ifthenelse{\boolean{bajasusua}}
{
\vspace{-25pt}
\begin{longtblr}[
	label = none,
	entry = none,
	]{
		width = \linewidth,
		colspec = {Q[935]},
		cells = {c},
                     row{1}={dorado},
		%hlines,
		%vlines,
	}
	\textbf{BAJAS} \\Sección D.- Accesos de usuarios
\end{longtblr}
%%%PARA PODER ALMACENAR TODAS LAS FILAS DEL DOCUMENTOS CSV ES NECESARIO TODO ESTO
\vspace{-30pt}
 \begin{longtblr}[
 label = none,
 entry = none,
 ]{
  width = \linewidth,
  colspec = {Q[2]Q[10]Q[12]Q[10]Q[8]Q[15]Q[5]Q[5]},                     
  hlines,
  vlines,
                     cells={c},
 }
\textbf{No°}&\textbf {Nombre(s) de \\ usuario(s)\\zona(s)\\origen}&\textbf{Segmento(s)/IP(s) \\origen}&
\textbf{Sistema asociado a \\ anfitrión(es) destino}&
\textbf{Rol de \\anfitrión(es) \\destino}&\textbf{IP/NAT anfitrión(es) \\destino/NAT}&
\textbf{Protocolo\\ TCP o UDP}&\textbf{Puertos}
\end{longtblr}
%%%%USUARIO
\DTLforeach*{BAJASUSUA}{\No=N, \NombreUsua=SO, \IPOri=IPO, \SistemaDes=SD, \FuncionDes=FRD, \IPDes=IPD, \Protocolo=PRO, \Puertos=PUER}
{
\vspace{-37pt}
 \begin{longtblr}[
 label = none,
 entry = none,
 ]{
  width = \linewidth,
  colspec ={Q[2]Q[10]Q[12]Q[10]Q[8]Q[15]Q[5]Q[5]},                     
  hlines,
 vlines,
                     cells={c},
 }
\No&\NombreUsua&\IPOri&\SistemaDes& \FuncionDes&\IPDes&\Protocolo& \Puertos
\end{longtblr}
}
}
}%%FIN DE USUARIOS


%%%%%JUSTIFICACION
\vspace{-20pt}
\begin{longtblr}[
	label = none,
	entry = none,
	]{
		width = \linewidth,
		colspec = {Q[935]},
		cells = {c},
                     row{1}={dorado},
		row{2} = {},
		hlines,
		vlines,
	}
	\textbf{Justificación} \\ \JUSTIFICA
\end{longtblr}

%\end{landscape}
\clearpage
%%%..... POLÍTICAS
\textbf{Políticas del servicio}
{\begin{itemize}
	        \item El solicitante deberá presentar este formato adjuntando el Memorando o Atenta nota y número de ticket de Mesa de ayuda asociado, sin los cuales no se podrá atender su solicitud.  
	        \item El formato deberá estar debidamente llenado y contener toda la información requerida facilitando la aplicación expedita de las configuraciones solicitadas. Es responsabilidad del solicitante recabar la información con los Administradores de los sistemas o Áreas involucrados (Subgerencia de Sistemas y/o Subgerencia de Internet e Intranet).
           \item El solicitante deberá proporcionar la dirección IP física y si utiliza dirección NAT agregarlo por cada servidor involucrado, no siendo responsabilidad de la SSTTS si al establecer la comunicación no conecta por falta de este dato. De no proporcionarse la dirección IP NAT correcta, las reglas de cortafuegos se configurarán por defecto con la dirección IP del adaptador de red de los servidores y corresponderá al Administrador del sistema aplicar los cambios en el sistema o servidores para establecer la comunicación.
           \item El solicitante tiene la obligación de indicar si se trata de un traslado de permisos de una dirección IP a otra.
           \item Para el traslado de permisos de una dirección IP a otra, se deberá llenar la sección de BAJAS con los permisos de la dirección IP anterior que se darán de baja, además de llenar la sección de ALTAS / CAMBIOS con los permisos que se requieren trasladar.
           \item Si el solicitante NO indica que se trata de un traslado de permisos, éste será responsable de cualquier acceso no autorizado que se derive de los permisos de la dirección IP anterior debido a no haber solicitado la baja correspondiente de dichos accesos.
           \item Es responsabilidad de los administradores de cada sistema llevar un control de las direcciones IP’s con acceso al sistema que administra.
           \item Es responsabilidad de los administradores de cada sistema documentar el control de accesos en la Memoria Técnica. Se sugiere actualizarla periódicamente cada 2, 3, 4 o 6 meses.
           \item Es responsabilidad de los administradores de cada servidor llevar un control de las direcciones IP’s con acceso al servidor que administra.
           \item La solicitud para cambios en la infraestructura de seguridad (RFCs) será solicitada únicamente por los Administradores de cada sistema una vez que se apliquen los permisos de acceso en el propio sistema y aperturen los accesos en los cortafuegos locales de los servidores involucrados.
           \item Es responsabilidad de los Administradores documentar el control de accesos de cada servidor en su respectiva Memoria Técnica, actualizando la misma periódicamente cada 2, 3, 4 0 6 meses.
           \item Al firmar el solicitante se da por enterado de las políticas del servicio y acepta la responsabilidad de cualquier materialización de los riesgos derivados de las aperturas de comunicaciones asociadas al presente control de cambios.
\end{itemize}
\newboolean{enlacebool}
\setboolean{enlacebool}{\ENLACEBOOL}

\newboolean{solibool}
\setboolean{solibool}{\SOLIBOOL}

\newboolean{enlacesolibool}
\setboolean{enlacesolibool}{\ENLACESOLIBOOL}

%%AMBOS
\ifthenelse{\boolean{enlacesolibool}}{
         
           %\vspace{-7pt}
	\begin{longtblr}[
	label = none,
	entry = none,
	]{
			width = \linewidth,
			colspec = {Q[200]Q[200]Q[200]},
			%row{1} = {gris},
			row{1} = {c},
                                row{2}={c},
			%cell{1}{1} = {c=2}{0.942\linewidth},
                                cell{2}{1} = {c=2}{0.942\linewidth},
			%hlines,
			%vlines,
		}
		%Fecha: \today                      &                                                                                                       \\
		{Gerente o director local que autoriza\\[1cm] \rule{6cm}{0.5mm}\\\NOMBREJEFE\\\PUESTOJEFE} & 
                   {Usuario que acepta.\\[1cm] \rule{6cm}{0.5mm}\\\NOMS\\\PUESTOS} &
                   {Enlace Informático\\[1cm] \rule{6cm}{0.5mm}\\\NOMEI}
	\end{longtblr}
}

%%SOLI
          \ifthenelse{\boolean{solibool}}{
           %\vspace{-7pt}
	\begin{longtblr}[
	label = none,
	entry = none,
	]{
			width = \linewidth,
			colspec = {Q[200]Q[200]Q[200]},
			%row{1} = {gris},
			row{1} = {c},
                                row{2}={c},
			%cell{1}{1} = {c=2}{0.942\linewidth},
                                cell{2}{1} = {c=2}{0.942\linewidth},
			%hlines,
			%vlines,
		}
		%Fecha: \today                      &                                                                                                       \\
		{Gerente o director local que autoriza\\[1cm] \rule{6cm}{0.5mm}\\\NOMBREJEFE\\\PUESTOJEFE} & 
                   {Usuario que acepta.\\[1cm] \rule{6cm}{0.5mm}\\\NOMS\\\PUESTOS} 
                  % {Enlace Informático\\[1cm] \rule{6cm}{0.5mm}\\\NOMEI}
	\end{longtblr}
}

\ifthenelse{\boolean{enlacebool}}{         
           %\vspace{-7pt}
	\begin{longtblr}[
	label = none,
	entry = none,
	]{
			width = \linewidth,
			colspec = {Q[200]Q[200]Q[200]},
			%row{1} = {gris},
			row{1} = {c},
                                row{2}={c},
			%cell{1}{1} = {c=2}{0.942\linewidth},
                                cell{2}{1} = {c=2}{0.942\linewidth},
			%hlines,
			%vlines,
		}
		%Fecha: \today                      &                                                                                                       \\
		{Gerente o director local que autoriza\\[1cm] \rule{6cm}{0.5mm}\\\NOMBREJEFE\\\PUESTOJEFE} & 
                  % {Usuario que acepta.\\[1cm] \rule{6cm}{0.5mm}\\\NOMS\\\PUESTOS} &
                   {Enlace Informático\\[1cm] \rule{6cm}{0.5mm}\\\NOMEI}
	\end{longtblr}
}

\end{document}